\section{Limitations}
\label{sec:limits}

\paragraph{Alias divergence is invisible, and that is a property of the medium.}
When two writers mint different names for one entity, the graphs disagree about
the world while agreeing about every triple. No triple-level operator can see
it, so a recall of $1.000$ is not achievable by anything in Table~\ref{tab:armA}.
We left the class in the benchmark deliberately: it is the most common real
divergence defect in the recorded corpus, and excluding it would have produced a
shape-aware arm scoring $1.000$ on both precision and recall as an artifact of
what we chose to measure. Arm D tests a separate pairwise entity-resolution stage, using labels and
historical neighborhoods to propose identity decisions. Its held-out run recovered
$11/88$ aliases ($11/55$ with eligible evidence) and proposed one false identity.
Semantic recall was $1/45$ ($1/19$ eligible), compared with label matching's
$3/45$ ($3/19$ eligible) and zero false proposals. The preregistered baseline-beating
condition fails on this stratum, a negative finding for Jev. Candidate generation
is a separate requirement whose recall bounds end-to-end resolution, and we
measured it with free signals only. Each pair's left endpoint ranks every entity
asserting one of its types before the repair, by the larger of normalized label
similarity and MiniLM cosine over the same cleaned pre-repair evidence, with ties
counted against the true partner. Recall at $k = 1, 5, 10, 50$ is $38$, $54$,
$56$ and $57$ of $93$ (id-form $23, 32, 34, 35$ of $46$; semantic
$15, 22, 22, 22$ of $47$). The ceiling is the type block, not the ranking:
$30$ partners cannot be ranked at all, $26$ of them because they never enter the
anchor's typed pool ($15$ had no pre-repair facts and $11$ carried disjoint
types), and $25$ of the $30$ are semantic. Among the $63$ rankable partners,
$56$ fall in the top ten, and cosine alone puts more partners first ($43$) than
the combined score ($38$). A resolver fed by this blocker would lose about a
third of all aliases, and over half of the semantic class, before pairwise
resolution begins. Historical evidence gaps, the semantic negative
shortfall, easy reviewed negatives, and shared endpoints limit generalization.
No actual identities were merged, so the study measures proposed errors rather
than downstream damage or the economic cost of accepting them.

\paragraph{The $939$ is a counterfactual.} Those doublings arose from sequential
re-posts through a single write path, not from a share, a divergence and a merge.
The supported claim is narrow: had the same edits arrived as a divergence and
been reconciled by this operator, each would have been surfaced as a decision
instead of silently doubling. It is \emph{not} a claim that the operator was
running and prevented $1{,}784$ doublings.

\paragraph{Comment bodies in the merge replay corpus are synthetic.} Real values are prose
about a deployment. In arms A and B, only the multiplicity is real --- which is all the merge
reasons about, but it does mean the corpus cannot support any claim about value
content.

\paragraph{A conflict blocks the whole merge, not the conflicting slot.} The
operator returns \texttt{conflicts} with nothing asserted. On a slice like the
$939$-subject one, a single reconnect writes nothing until every contended slot
is resolved. This is the conservative choice and it is a real usability cost;
partial application with a held-back conflict set is future work.

\paragraph{Ground graphs only.} The benchmark's graphs contain no blank nodes,
so canonicalization there reduces to sorting. RDFC-1.0 is the substrate we build
on and is not exercised by this arm; nothing here is evidence about blank-node
canonicalization.

\paragraph{Timing is indicative.} One machine, one process, single runs. The
line-merge arms include process spawn for \texttt{git merge-file}, of order
$5$\,ms, which dominates their timing at every size measured; those columns are
not comparable with the in-process arms and we do not report them as if they
were.

\paragraph{Slot granularity is an assumption.} A conflict is reported over a
$(subject, predicate)$ slot because that is the unit \texttt{sh:maxCount} is
declared over. An operator working at another granularity would need a different
scorer, and the comparison in Table~\ref{tab:armA} would not transfer unchanged.

\paragraph{One deployment.} The recorded corpus comes from a single production
store. Its class distribution --- alias-mint and functional-value doubling
dominating --- may not generalize, and we make no claim about divergence rates
in other systems from it.

\paragraph{Private evidence limits inference reproducibility.}
Arm D publishes the item labels, frozen free-baseline scores, prompt, numeric
responses, attempt ledger, and offline scorer. Its structural item views replace
private strings with random tokens and withhold descriptions; selected evidence
examples are scrubbed editorial paraphrases. These artifacts reproduce scoring,
not the original model requests or embeddings. Reviewer inspection used the
preserved private evidence. A reader cannot independently regenerate the model
outputs from the public views, and a commercial model may change separately.
Timeouts may be billed without usage responses, so list-price subtotals do not
establish total account spend. First-attempt and eventual-item availability
must be reported separately when retries are allowed.
